\section{Introduction}
Throughout, $\NN=\{0,1,2,\ldots\}$. A set $S\subseteq\NN^n$ is
\emph{recursively enumerable} if a deterministic computation accepts exactly
its members; it may run forever on other inputs. A \emph{Diophantine
representation} of $S$ is a polynomial
\[
 P\in\ZZ[A_1,\ldots,A_n,W_1,\ldots,W_m]
\]
such that $a\in S$ if and only if $P(a,w)=0$ for some $w\in\NN^m$.
The representation is \emph{single-fold} if this complete tuple $w$ is unique
whenever it exists. It is \emph{finite-fold} if the set of such tuples is finite
for every input. These definitions count all auxiliary variables, including
those used to express subsidiary arithmetic operations.

\begin{theorem}\label{thm:main}
For every integer $n\ge1$ and every recursively enumerable set
$S\subseteq\NN^n$, there are an integer $m\ge1$ and a polynomial
$P\in\ZZ[A_1,\ldots,A_n,W_1,\ldots,W_m]$ such that, for every $a\in\NN^n$,
\[
 \#\{w\in\NN^m:P(a,w)=0\}=
 \begin{cases}1,&a\in S,\\0,&a\notin S.\end{cases}
\]
\end{theorem}

The Davis--Putnam--Robinson theorem established exponential Diophantine
representations of recursively enumerable relations: the defining expressions
may use exponentiation as well as addition and multiplication
\cite[Section~1]{dpr1961}. Robinson had shown how suitable growth conditions
yield an existential definition of exponentiation
\cite[Section~3]{robinson1952}. Matiyasevich's 1970 construction of an
exponentially growing Diophantine relation completed this route to polynomial
representations \cite[Sections~1 and~3]{mati1970}. The resulting representation
theorem implies the negative solution of Hilbert's tenth problem, since a
decision procedure for polynomial solvability would decide membership in every
recursively enumerable set. It does not bound the number of witnesses.

Matiyasevich obtained a different strengthening in 1974: every recursively
enumerable relation has an exponential representation with a unique complete
auxiliary tuple \cite[Section~2, Main Theorem]{mati1974}. The same paper
already formulates the remaining problem of polynomial representations with
this uniqueness and reduces it to single-fold binary exponentiation
\cite[Section~1]{mati1974}. This is the single-fold conjecture; allowing
finitely many witnesses gives the finite-fold conjecture \cite{mati2010}.
Theorem~\ref{thm:main} proves the single-fold assertion over $\NN$.
The witness domain matters: Miller and Shlapentokh obtain single-fold
definitions of computably enumerable sets of rational integers using
witnesses in characteristic-zero polynomial rings, and finite-fold definitions
of computably enumerable relations in polynomial rings over rings of
integers of totally real number fields \cite{millershlapentokh2022}.

Our arithmetic construction develops the classical strategy of recovering
powers from Pell solutions and controlling their indices through growth
bounds. Cantone, Cuzziol, and Omodeo give a recent account, including a
paired-Pell representation and the problem of replacing its exponential
cutoff \cite[Section~6, Lemma~6.1 and Theorem~5; Section~7]{cco2024}.
We supply a cutoff using rational points on a rank-one elliptic curve.
Poonen uses rank-one curves and denominator growth to obtain Diophantine
definitions between rings of integers
\cite[Theorem~1 and Section~2.3, Lemma~8]{poonen2002}.
Here the N\'eron--Tate height compares the full logarithmic x-coordinate
height with a quadratic function of an integer index, with uniformly bounded
additive error \cite[Chapter~VIII, Proposition~9.1 and Theorem~9.3]{silverman}.
The bounded additive error makes the exact index selection below possible.

The computational construction uses the bit-masking method of Jones and
Matiyasevich \cite{jm1984,jm1991}. Large-base encodings turn a finite
register-machine history into finitely many integers, and Lucas's binomial
congruence tests the permitted binary digits. Determining the base from the
input and computation time already appears in Chaitin's account
\cite[Sections~2.3 and~2.6]{chaitin1987}. We use a clock and a dyadic
interval to select one base for the entire run. Throughout both constructions,
uniqueness must include reduced rational coordinates, arithmetic quotients
and slacks, and the variables of inactive cases. The proof keeps this
complete-witness requirement through every substitution.

The arithmetic constants in the proof are fixed once for all sets $S$ and
all inputs. They include rational points and integers whose existence is
established below, but whose numerical values are not specified. The theorem
therefore gives existence of an integer polynomial without displaying its
coefficients or a procedure for numerically choosing those constants.

\begin{corollary}[Undecidability with an at-most-one-solution promise]
\label{cor:unique-h10}
There exist integers $m,D\ge1$ for which no algorithm has the following
property: on every input polynomial $F\in\ZZ[W_1,\ldots,W_m]$ of total
degree at most $D$ satisfying the promise
\[
 \#\{w\in\NN^m:F(w)=0\}\le1,
\]
the algorithm halts and correctly decides whether this set is nonempty.
\end{corollary}
